\chapter{重力理論について}
kkk
\section{ニュートン力学}
\cite{1971430859808286867}
ニュートン力学はニュートン(Isaac Newton)が1687年にプリンキピア(Principia)で発表した古典力学の枠組みである．当時，2つの天体を結びつける力の正体を多くの研究者が探っていた．その中でニュートンはすべての物体同士は互いに引っ張り合う力が働くとする万有引力を新たに考えた．さらに，地球上での物体の落下現象と天体間の運動は同一の力が起因していると説明し，万有引力の法則を提唱した．ニュートン力学は運動の3法則と万有引力の法則から成る理論形式で，ニュートン力学において重力とは，万有引力と地球の自転による遠心力の合力である．しかし，古典力学では原子レベルのミクロな世界での物理法則や光速に近い速度の物体の運動を扱うことができなかった．それらを解決するために提唱されたのが量子論と相対論である．その点でニュートン力学は一般性があるとは言えず，重力の捉え方も見直す必要があった．
\section{一般相対性理論}
\cite{1971149384824946599,1020000782474158465}
一般相対性理論(GR:General Relativity)は1915年にアインシュタイン(Albert Einstein)が提唱した理論である．これは1905年に提唱した特殊相対性理論がもとになっている．当時，光は未知の媒質(エーテル)を振動しながら伝わると考えられていた．しかし，マイケルソン・モーリーの実験より光の速度は方向による違いは見られず，エーテルの存在は確認できなかった．この実験を皮切りに，アインシュタインは光速度不変の原理と相対性原理を見出し，特殊相対性理論を発表した．
ニュートン力学では重力は無限の速さで伝わることになるが，特殊相対論では光速を超えて情報を伝えることはできない．
\section{場の量子論}
\cite{1971712334789710479,1971993809681112242}
